\documentclass[12pt,a4paper]{article}
\usepackage[margin=1in]{geometry}
\usepackage{hyperref,amsmath,amssymb,booktabs,longtable,fancyhdr,array,microtype}
\hypersetup{colorlinks=true,linkcolor=blue!70!black,urlcolor=blue!70!black}
\pagestyle{fancy}\fancyhf{}
\fancyhead[L]{Agentic AI Bootcamp}
\fancyhead[R]{sujeet.banerjee@gmail.com\\sujeet.banerjee.professional@gmail.com}
\fancyfoot[L]{Langfuse SaaS Setup Guide}
\fancyfoot[R]{\thepage}
\renewcommand{\headrulewidth}{0.4pt}\renewcommand{\footrulewidth}{0.4pt}
\title{\textbf{Langfuse SaaS Setup \& LiteLLM Integration}\\\large Step-by-Step Guide for Agentic Workflows}
\author{Detailed Notes -- Agentic AI Bootcamp}\date{}
\begin{document}\maketitle\tableofcontents\newpage

\section{Introduction}
Observability is a critical component of any modern AI architecture. When orchestrating agentic workflows using platforms like n8n and routing LLM calls through a proxy model, a dedicated telemetry platform is required to trace execution, monitor latency, and analyze token costs. Langfuse provides an observability platform designed specifically for LLM applications.

This guide outlines the process of creating a free Langfuse Cloud account and properly integrating it with a LiteLLM proxy deployed in a GitHub Codespace.

\section{CONFIGURING LANGFUSE}

\subsection{Step 1: Creating a Free Langfuse Account}
Langfuse offers a fully managed cloud tier with a generous free plan suitable for development and bootcamp environments.

\begin{enumerate}
    \item Navigate to the Langfuse Cloud portal. You must choose a data region:
    \begin{itemize}
        \item \textbf{US Region:} \url{https://us.cloud.langfuse.com}
        \item \textbf{EU Region:} \url{https://cloud.langfuse.com}
    \end{itemize}
    \item Select \textbf{Sign Up} and authenticate using your preferred method (GitHub, Google, or Email/Password).
    \item Upon initial login, you will be prompted to create an \textbf{Organization} and a \textbf{Project}. Give your project a descriptive name, such as \textit{``Agentic Crew Bootcamp''}.
\end{enumerate}

\subsection{Step 2: Generating API Credentials}
To allow the local environment to transmit telemetry data, you must generate API keys.

\begin{enumerate}
    \item In your Langfuse dashboard, navigate to \textbf{Settings} (the gear icon) on the left navigation menu.
    \item Select \textbf{API Keys}.
    \item Click \textbf{Create New API Key}.
    \item A modal will appear displaying three critical pieces of information. \textbf{Copy these immediately}, as the Secret Key will not be shown again:
    \begin{itemize}
        \item \texttt{Public Key} (Starts with \texttt{pk-lf-...})
        \item \texttt{Secret Key} (Starts with \texttt{sk-lf-...})
        \item \texttt{Host URL} (Matches the region you signed up in, e.g., \url{https://jp.cloud.langfuse.com})
    \end{itemize}
\end{enumerate}

\section{CONSUMING LANGFUSE CREDENTIALS}

\subsection{Ex 1: Configuring the Github \textit{Codespace} Environment}
LiteLLM utilizes the Langfuse Python SDK internally. The SDK requires specific environment variables to authenticate and route payloads correctly.

In your GitHub Codespace, locate your environment configuration file (typically \texttt{.env}) and append the following variables:

\begin{verbatim}
LANGFUSE_PUBLIC_KEY="pk-lf-YOUR_PUBLIC_KEY"
LANGFUSE_SECRET_KEY="sk-lf-YOUR_SECRET_KEY"
LANGFUSE_HOST="https://jp.cloud.langfuse.com" 
\end{verbatim}

\begin{quote}
\textbf{Crucial Architectural Note:} The host environment variable \textit{must} be named \texttt{LANGFUSE\_HOST}. Using alternatives like \texttt{LANGFUSE\_BASE\_URL} will cause the OpenTelemetry exporter to silently fall back to its default US endpoints. This results in \texttt{401 Invalid credentials} errors during the asynchronous batch export if your keys belong to a different geographic region.
\end{quote}

\subsubsection{Updating the LiteLLM Proxy Configuration}
Ensure your \texttt{config.yaml} explicitly enables Langfuse in the callbacks array.

\begin{verbatim}
# LiteLLM Model configs
...
...

# LangFuse OTEL config
litellm_settings:
  callbacks:
    - langfuse_otel
\end{verbatim}

Restart your Docker container to apply the new environment variables and configuration:
\begin{verbatim}
docker restart litellm_proxy
\end{verbatim}

\subsection{EX 2: Codespace settings - secrets injection}
An alternative approach to using .env, is to use codespace settings to add the LangFuse secrets here: https://github.com/settings/codespaces

For more details, see the reference documents:
\begin{enumerate}
    \item https://dev.to/github/securely-store-environment-variables-with-github-codespaces-3dgc
    \item https://docs.github.com/en/codespaces/developing-in-a-codespace/persisting-environment-variables-and-temporary-files 
\end{enumerate}


% \subsection*{Step 5: End-to-End Verification}
% \begin{enumerate}
%     \item Trigger an AI workflow from your n8n instance that invokes the chat model endpoint.
%     \item Verify the container logs to ensure no asynchronous export errors (such as 401 or 403) occur.
%     \item Open the Langfuse SaaS dashboard and navigate to the \textbf{Traces} or \textbf{Generations} tab.
%     \item You should now see detailed telemetry for your request, including the model used (e.g., \texttt{llama-3.3-70b-versatile}), token usage, latency, and inferred cost.
% \end{enumerate}

\vfill
\begin{center}\small
\textcopyright\ 2026 Sujeet Banerjee. All rights reserved.\\
Agentic AI Bootcamp -- Educational material.
\end{center}
\end{document}
